\chapter*{Nomenclature}
\addcontentsline{toc}{chapter}{Nomenclature}
\markboth{Nomenclature}{Nomenclature}

\section*{General Conventions}
\begin{itemize}
    \item Scalars are lowercase or Greek letters (e.g., $m, t, \theta$).
    \item Vectors are bold lowercase (e.g., $\mathbf{x}, \mathbf{u}, \mathbf{q}$).
    \item Matrices are bold uppercase (e.g., $\mathbf{A}, \mathbf{B}, \mathbf{M}, \mathbf{J}$).
    \item Expected values and sets are denoted by blackboard bold or script (e.g., $\mathbb{E}, \mathcal{S}$).
\end{itemize}

These are default conventions; each chapter declares its model, frames, and
any local reuse of symbols. The mechanical entries below use $n$ independent,
regular local configuration coordinates and $\mathbf{v}=\dot{\mathbf{q}}$.
This convention is local: it does not identify every configuration manifold
with a Euclidean space or apply unchanged to redundant representations.

\section*{State and Kinematics}
\begin{description}
    \item[$\mathbf{q} \in \mathbb{R}^n$] Configuration coordinates in the declared local chart.
    \item[$\dot{\mathbf{q}} \in \mathbb{R}^n$] Coordinate rates; the generalized velocities under the local convention above.
    \item[$\ddot{\mathbf{q}} \in \mathbb{R}^n$] Second time derivatives of these coordinates, not necessarily Cartesian point accelerations.
    \item[$\mathbf{x} \in \mathbb{R}^{n_x}$] State-vector representation with $n_x$ components; commonly $\mathbf{x} = [\mathbf{q}^T, \dot{\mathbf{q}}^T]^T$ locally, with actuator or internal states added when needed. Redundant state representations also retain their constraints.
    \item[$\mathbf{u} \in \mathbb{R}^m$] Control input vector (e.g., joint torques, muscle activations).
    \item[$\boldsymbol{\tau} \in \mathbb{R}^n$] Generalized forces conjugate to $\mathbf{q}$, with power $\boldsymbol{\tau}^T\dot{\mathbf{q}}$; components need not all be joint torques.
    \item[$t \in \mathbb{R}$] Time.
\end{description}

More generally, write $\mathbf{q}\in\mathbb{R}^{n_q}$ for a configuration
representation and $\mathbf{v}\in\mathbb{R}^{n_v}$ for generalized velocities,
with $\dot{\mathbf{q}}=\mathbf{N}(\mathbf{q})\mathbf{v}$ on the admissible domain.
A unit quaternion has four components with a norm constraint, while angular
velocity has three components in a declared frame. Configuration rates and
generalized velocities therefore need not have the same dimension. See
\href{https://underactuated.mit.edu/multibody.html}{Tedrake's multibody notes}
for this representation convention. Constraints and any time-dependent
kinematics must be included in the chosen model.

\section*{Inertia and Dynamics}
\begin{description}
    \item[$\mathbf{M}(\mathbf{q})$] Mass (inertia) matrix. It is symmetric positive definite for independent regular velocities with strictly positive modeled kinetic energy for every nonzero velocity. Constrained, redundant, massless, or singular representations require their own treatment.
    \item[$\mathbf{C}(\mathbf{q}, \dot{\mathbf{q}})$] Coriolis/centrifugal matrix; $\mathbf{C}\dot{\mathbf{q}}$ is the generalized velocity-product force in the local mechanical equation. The matrix factorization is not unique.
    \item[$\mathbf{g}(\mathbf{q})$] Generalized gravity term; its sign follows the side of the mechanical equation on which the chapter places it.
    \item[$\mathbf{J}(\mathbf{q})$] Coordinate Jacobian of a time-independent task map $\mathbf{p}(\mathbf{q})$, so $\dot{\mathbf{p}}=\mathbf{J}\dot{\mathbf{q}}$ in the stated coordinates and frames.
\end{description}

With $\dot{\mathbf{q}}=\mathbf{N}\mathbf{v}$, the corresponding velocity Jacobian is
$\mathbf{J}_v=\mathbf{J}_q\mathbf{N}$. Force components must change with the velocity convention:
$\boldsymbol{\tau}_v=\mathbf{N}^T\boldsymbol{\tau}_q$ gives
$\boldsymbol{\tau}_v^T\mathbf{v}=\boldsymbol{\tau}_q^T\dot{\mathbf{q}}$.
Thus the same club-point motion and mechanical power must result from
consistently paired kinematic and force representations.

\section*{Control and Optimization}
For the indexed entries, declare a discrete update
$\mathbf{x}_{k+1}=F_k(\mathbf{x}_k,\mathbf{u}_k)$ and a nominal pair
$(\bar{\mathbf{x}}_k,\bar{\mathbf{u}}_k)$. Its derivatives are evaluated at
that pair. These are derivatives of the full update, including the chosen
sampling or numerical integration rule, not of zero-input drift alone.
\begin{description}
    \item[$\mathbf{A}_k = D_x F_k$] State Jacobian of the declared discrete update at the nominal pair.
    \item[$\mathbf{B}_k = D_u F_k$] Input Jacobian of that update at the same pair.
    \item[$V(\mathbf{x})$ or $\mathcal{V}$] Value function or Lyapunov function.
    \item[$\mathbf{Q}$] State penalty matrix in LQR/DDP.
    \item[$\mathbf{R}$] Control penalty matrix in LQR/DDP.
    \item[$\mathbf{K}$] Feedback gain matrix ($\delta\mathbf{u} = -\mathbf{K}\delta\mathbf{x}$).
    \item[$\boldsymbol{\gamma}$] A curve with a declared parameter; a timed state trajectory $\boldsymbol{\gamma}(t)$ fixes an initial state and input law. A solution map $\phi(t;t_0,\mathbf{x}_0)$ instead retains dependence on initial time and state, for a fixed evolution law on its domain of existence.
\end{description}

Continuous-time chapters define their Jacobians from the full differential
equation. Some use $f(\mathbf{x},\mathbf{u},t)$ for that full right-hand side;
the autonomous $f(\mathbf{x})$ below denotes only the declared zero-input drift.
Neither convention makes a discrete input Jacobian equal to the continuous
one without accounting for the sampling interval and evolution.

\section*{Affine Control and Counterfactuals}
\begin{description}
    \item[$f(\mathbf{x})$] Drift vector field: complete autonomous evolution of the declared effective plant when $\mathbf{u} = \mathbf{0}$.
    \item[$G(\mathbf{x})$] Input (control) matrix field mapping control inputs to state derivatives.
    \item[$\dot{\mathbf{x}} = f(\mathbf{x}) + G(\mathbf{x})\mathbf{u}$] Control-affine system equations of motion.
    \item[$\mathbf{x}^{\mathrm{ZTCF}}(t)$] Forward Zero-Torque Counterfactual trajectory under $\mathbf{u}(t) \equiv \mathbf{0}$.
    \item[$\mathbf{a}_{\mathrm{ZVCF}}(\mathbf{q}, \mathbf{z})$] Instantaneous Zero-Velocity Counterfactual acceleration with velocity and declared control set to zero.
    \item[$\mathrm{DCR}_{W,\mathcal{U}}(\mathbf{x})$] Drift-Control Ratio in a declared acceleration or task-projected space:
    \[
      \mathrm{DCR}_{W,\mathcal{U}}(\mathbf{x})
      =\frac{\|W \mathbf{a}_{\mathrm{drift}}\|_2}{\sup_{\mathbf{u}\in\mathcal{U}}\|W B_a \mathbf{u}\|_2 + \varepsilon}.
    \]
\end{description}

\section*{Biomechanics and Motor Control}
\begin{description}
    \item[{$a_i \in [0, 1]$}] Muscle activation level.
    \item[$\ell_i^M$] Muscle fiber length.
    \item[$v_i^M$] Muscle fiber contraction velocity.
    \item[$\ell_i^{MT}$] Musculotendon unit length.
    \item[$\mathbf{F}_{\text{muscle}}$] Force generated by muscles.
    \item[$\mathbf{W}$] Muscle synergy matrix (weights).
    \item[$\mathbf{c}(t)$] Muscle synergy activation vector over time.
\end{description}

\section*{Geometry and Manifolds}
\begin{description}
    \item[$SE(3)$] Special Euclidean group of rigid body motions.
    \item[$SO(3)$] Special Orthogonal group of 3D rotations.
    \item[$\mathfrak{se}(3)$] Lie algebra of $SE(3)$, space of spatial twists (velocities).
    \item[$\mathcal{M}$] A smooth manifold.
    \item[$T_{\mathbf{x}}\mathcal{M}$] Tangent space at point $\mathbf{x}$.
\end{description}
